%=====================================
%   osmh.tex
%   osmh
%   toshi2019
%=====================================

% -----------------------
% preamble
% -----------------------
% ここから本文 (\begin{document}) までの
% ソースコードに変更を加えた場合は
% 編集者まで連絡してください. 
% Don't change preamble code yourself. 
% If you add something
% (usepackage, newtheorem, newcommand, renewcommand),
% please tell it 
% to the editor of institutional paper of RUMS.

% ------------------------
% documentclass
% ------------------------
\documentclass[leqno, 10pt, a4paper, dvipdfmx]{jsarticle}

% ------------------------
% usepackage
% ------------------------
\usepackage{algorithm}
\usepackage{algorithmic}
\usepackage{amscd}
\usepackage{amsfonts}
\usepackage{amsmath}
\usepackage[psamsfonts]{amssymb}
\usepackage{amsthm}
\usepackage{ascmac}
\usepackage{bm}
\usepackage{color}
\usepackage{enumerate}
\usepackage{fancybox}
\usepackage[stable]{footmisc}
\usepackage{graphicx}
\usepackage{listings}
\usepackage{mathrsfs}
\usepackage{mathtools}
\usepackage{otf}
\usepackage{pifont}
\usepackage{proof}
\usepackage{subfigure}
\usepackage{tikz}
\usepackage{verbatim}
\usepackage[all]{xy}

\usetikzlibrary{cd}
\usetikzlibrary{arrows.meta}


% ================================
% パッケージを追加する場合のスペース 
\usepackage[dvipdfmx]{hyperref}
%\usepackage{xcolor}
\definecolor{darkgreen}{rgb}{0,0.45,0} 
\definecolor{darkred}{rgb}{0.75,0,0}
\definecolor{darkblue}{rgb}{0,0,0.6} 
\hypersetup{
    colorlinks=true,
    %citecolor=darkgreen,
    %linkcolor=darkred,
    %urlcolor=darkblue,
    citecolor=black,
    linkcolor=black,
    urlcolor=black,
}
\usepackage{pxjahyper}

%\usepackage[mathscr]{euscript} % mathscr のフォントを変更

%\usepackage{mmasym}

%=================================


% --------------------------
% theoremstyle
% --------------------------
\theoremstyle{definition}

% --------------------------
% newtheoem
% --------------------------

% 日本語で定理, 命題, 証明などを番号付きで用いるためのコマンドです. 
% If you want to use theorem environment in Japanece, 
% you can use these code. 
% Attention!
% All theorem enivironment numbers depend on 
% only section numbers.
\newtheorem{Axiom}{公理}%[section]
\newtheorem{Definition}[Axiom]{定義}
\newtheorem{Theorem}[Axiom]{定理}
\newtheorem{Proposition}[Axiom]{命題}
\newtheorem{Lemma}[Axiom]{補題}
\newtheorem{Corollary}[Axiom]{系}
\newtheorem{Example}[Axiom]{例}
\newtheorem{Claim}[Axiom]{主張}
\newtheorem{Property}[Axiom]{性質}
\newtheorem{Attention}[Axiom]{注意}
\newtheorem{Question}[Axiom]{問}
\newtheorem{Problem}[Axiom]{問題}
\newtheorem{Consideration}[Axiom]{考察}
\newtheorem{Alert}[Axiom]{警告}
\newtheorem{Fact}[Axiom]{事実}


% 日本語で定理, 命題, 証明などを番号なしで用いるためのコマンドです. 
% If you want to use theorem environment with no number in Japanese, You can use these code.
\newtheorem*{Axiom*}{公理}
\newtheorem*{Definition*}{定義}
\newtheorem*{Theorem*}{定理}
\newtheorem*{Proposition*}{命題}
\newtheorem*{Lemma*}{補題}
\newtheorem*{Example*}{例}
\newtheorem*{Corollary*}{系}
\newtheorem*{Claim*}{主張}
\newtheorem*{Property*}{性質}
\newtheorem*{Attention*}{注意}
\newtheorem*{Question*}{問}
\newtheorem*{Problem*}{問題}
\newtheorem*{Consideration*}{考察}
\newtheorem*{Alert*}{警告}
\newtheorem{Fact*}{事実}


% 英語で定理, 命題, 証明などを番号付きで用いるためのコマンドです. 
%　If you want to use theorem environment in English, You can use these code.
%all theorem enivironment number depend on only section number.
\newtheorem{Axiom+}{Axiom}%[section]
\newtheorem{Definition+}[Axiom+]{Definition}
\newtheorem{Theorem+}[Axiom+]{Theorem}
\newtheorem{Proposition+}[Axiom+]{Proposition}
\newtheorem{Lemma+}[Axiom+]{Lemma}
\newtheorem{Example+}[Axiom+]{Example}
\newtheorem{Corollary+}[Axiom+]{Corollary}
\newtheorem{Claim+}[Axiom+]{Claim}
\newtheorem{Property+}[Axiom+]{Property}
\newtheorem{Attention+}[Axiom+]{Attention}
\newtheorem{Question+}[Axiom+]{Question}
\newtheorem{Problem+}[Axiom+]{Problem}
\newtheorem{Consideration+}[Axiom+]{Consideration}
\newtheorem{Alert+}{Alert}
\newtheorem{Fact+}[Axiom+]{Fact}
\newtheorem{Remark+}[Axiom+]{Remark}

% ----------------------------
% commmand
% ----------------------------
% 執筆に便利なコマンド集です. 
% コマンドを追加する場合は下のスペースへ. 

% 集合の記号 (黒板文字)
\newcommand{\NN}{\mathbb{N}}
\newcommand{\ZZ}{\mathbb{Z}}
\newcommand{\QQ}{\mathbb{Q}}
\newcommand{\RR}{\mathbb{R}}
\newcommand{\CC}{\mathbb{C}}
\newcommand{\PP}{\mathbb{P}}
\newcommand{\KK}{\mathbb{K}}


% 集合の記号 (太文字)
\newcommand{\nn}{\mathbf{N}}
\newcommand{\zz}{\mathbf{Z}}
\newcommand{\qq}{\mathbf{Q}}
\newcommand{\rr}{\mathbf{R}}
\newcommand{\cc}{\mathbf{C}}
\newcommand{\pp}{\mathbf{P}}
\newcommand{\kk}{\mathbf{k}}

% 特殊な写像の記号
\newcommand{\ev}{\mathop{\mathrm{ev}}\nolimits} % 値写像
\newcommand{\pr}{\mathop{\mathrm{pr}}\nolimits} % 射影

% スクリプト体にするコマンド
%   例えば　{\mcal C} のように用いる
\newcommand{\mcal}{\mathcal}

% 花文字にするコマンド 
%   例えば　{\h C} のように用いる
\newcommand{\h}{\mathscr}

% ヒルベルト空間などの記号
\newcommand{\F}{\mcal{F}}
\newcommand{\X}{\mcal{X}}
\newcommand{\Y}{\mcal{Y}}
\newcommand{\Hil}{\mcal{H}}
\newcommand{\RKHS}{\Hil_{k}}
\newcommand{\Loss}{\mcal{L}_{D}}
\newcommand{\MLsp}{(\X, \Y, D, \Hil, \Loss)}

% 偏微分作用素の記号
\newcommand{\p}{\partial}

% 角カッコの記号 (内積は下にマクロがあります)
\newcommand{\lan}{\langle}
\newcommand{\ran}{\rangle}



% 圏の記号など
\newcommand{\Set}{{\bf Set}}
\newcommand{\Vect}{{\bf Vect}}
\newcommand{\FDVect}{{\bf FDVect}}
%\newcommand{\Ring}{{\bf Ring}}
\newcommand{\Ab}{{\bf Ab}}
\newcommand{\Mod}{\mathop{\mathrm{Mod}}\nolimits}
\newcommand{\Modf}{\mathop{\mathrm{Mod}^\mathrm{f}}\nolimits}
\newcommand{\CGA}{{\bf CGA}}
\newcommand{\GVect}{{\bf GVect}}
\newcommand{\Lie}{{\bf Lie}}
\newcommand{\dLie}{{\bf Liec}}



% 射の集合など
\newcommand{\Map}{\mathop{\mathrm{Map}}\nolimits} % 写像の集合
\newcommand{\Hom}{\mathop{\mathrm{Hom}}\nolimits} % 射集合
\newcommand{\End}{\mathop{\mathrm{End}}\nolimits} % 自己準同型の集合
\newcommand{\Aut}{\mathop{\mathrm{Aut}}\nolimits} % 自己同型の集合
\newcommand{\Mor}{\mathop{\mathrm{Mor}}\nolimits} % 射集合
\newcommand{\Ker}{\mathop{\mathrm{Ker}}\nolimits} % 核
\newcommand{\Img}{\mathop{\mathrm{Im}}\nolimits} % 像
\newcommand{\Cok}{\mathop{\mathrm{Coker}}\nolimits} % 余核
\newcommand{\Cim}{\mathop{\mathrm{Coim}}\nolimits} % 余像

% その他便利なコマンド
\newcommand{\dip}{\displaystyle} % 本文中で数式モード
\newcommand{\e}{\varepsilon} % イプシロン
\newcommand{\dl}{\delta} % デルタ
\newcommand{\pphi}{\varphi} % ファイ
\newcommand{\ti}{\tilde} % チルダ
\newcommand{\pal}{\parallel} % 平行
\newcommand{\op}{{\rm op}} % 双対を取る記号
\newcommand{\lcm}{\mathop{\mathrm{lcm}}\nolimits} % 最小公倍数の記号
\newcommand{\Probsp}{(\Omega, \F, \P)} 
\newcommand{\argmax}{\mathop{\rm arg~max}\limits}
\newcommand{\argmin}{\mathop{\rm arg~min}\limits}





% ================================
% コマンドを追加する場合のスペース 
%\newcommand{\OO}{\mcal{O}}



\makeatletter
\renewenvironment{proof}[1][\proofname]{\par
  \pushQED{\qed}%
  \normalfont \topsep6\p@\@plus6\p@\relax
  \trivlist
  \item[\hskip\labelsep
%        \itshape
         \bfseries
%    #1\@addpunct{.}]\ignorespaces
    {#1}]\ignorespaces
}{%
  \popQED\endtrivlist\@endpefalse
}
\makeatother

\renewcommand{\proofname}{Proof.}



%\renewcommand\proofname{\bf 証明} % 証明
\numberwithin{equation}{section}
\newcommand{\cTop}{\textsf{Top}}
%\newcommand{\cOpen}{\textsf{Open}}
\newcommand{\Op}{\mathop{\textsf{Op}}\nolimits}
\newcommand{\Ob}{\mathop{\textrm{Ob}}\nolimits}
\newcommand{\id}{\mathop{\mathrm{id}}\nolimits}
\newcommand{\pt}{\mathop{\mathrm{pt}}\nolimits}
\newcommand{\res}{\mathop{\rho}\nolimits}
\newcommand{\A}{\mcal{A}}
\newcommand{\B}{\mcal{B}}
\newcommand{\C}{\mcal{C}}
\newcommand{\D}{\mcal{D}}
\newcommand{\E}{\mcal{E}}
\newcommand{\G}{\mcal{G}}
%\newcommand{\H}{\mcal{H}}
\newcommand{\I}{\mcal{I}}
\newcommand{\J}{\mcal{J}}
\newcommand{\OO}{\mcal{O}}
\newcommand{\Ring}{\mathop{\textsf{Ring}}\nolimits}
\newcommand{\cAb}{\mathop{\textsf{Ab}}\nolimits}
%\newcommand{\Ker}{\mathop{\mathrm{Ker}}\nolimits}
\newcommand{\im}{\mathop{\mathrm{Im}}\nolimits}
\newcommand{\Coker}{\mathop{\mathrm{Coker}}\nolimits}
\newcommand{\Coim}{\mathop{\mathrm{Coim}}\nolimits}
\newcommand{\rank}{\mathop{\mathrm{rank}}\nolimits}
\newcommand{\Ht}{\mathop{\mathrm{Ht}}\nolimits}
\newcommand{\supp}{\mathop{\mathrm{supp}}\nolimits}
\newcommand{\colim}{\mathop{\mathrm{colim}}}
\newcommand{\Tor}{\mathop{\mathrm{Tor}}\nolimits}

\newcommand{\cat}{\mathscr{C}}

%筆記体
\newcommand{\cA}{\mcal{A}}
\newcommand{\cB}{\mcal{B}}
\newcommand{\cC}{\mcal{C}}
\newcommand{\cD}{\mcal{D}}
\newcommand{\cE}{\mcal{E}}
\newcommand{\cF}{\mcal{F}}
\newcommand{\cG}{\mcal{G}}
\newcommand{\cH}{\mcal{H}}
\newcommand{\cI}{\mcal{I}}
\newcommand{\cJ}{\mcal{J}}
\newcommand{\cK}{\mcal{K}}
\newcommand{\cL}{\mcal{L}}
\newcommand{\cM}{\mcal{M}}
\newcommand{\cN}{\mcal{N}}
\newcommand{\cO}{\mcal{O}}
\newcommand{\cP}{\mcal{P}}
\newcommand{\cQ}{\mcal{Q}}
\newcommand{\cR}{\mcal{R}}
\newcommand{\cS}{\mcal{S}}
\newcommand{\cT}{\mcal{T}}
\newcommand{\cU}{\mcal{U}}
\newcommand{\cV}{\mcal{V}}
\newcommand{\cW}{\mcal{W}}
\newcommand{\cX}{\mcal{X}}
\newcommand{\cY}{\mcal{Y}}
\newcommand{\cZ}{\mcal{Z}}


\newcommand{\scA}{\mathscr{A}}
\newcommand{\scB}{\mathscr{B}}
\newcommand{\scC}{\mathscr{C}}
\newcommand{\scD}{\mathscr{D}}
\newcommand{\scE}{\mathscr{E}}
\newcommand{\scF}{\mathscr{F}}
\newcommand{\scN}{\mathscr{N}}
\newcommand{\scO}{\mathscr{O}}
\newcommand{\scV}{\mathscr{V}}
\newcommand{\scU}{\mathscr{U}}


\newcommand{\ibA}{\mathop{\text{\textit{\textbf{A}}}}}
\newcommand{\ibB}{\mathop{\text{\textit{\textbf{B}}}}}
\newcommand{\ibC}{\mathop{\text{\textit{\textbf{C}}}}}
\newcommand{\ibD}{\mathop{\text{\textit{\textbf{D}}}}}
\newcommand{\ibE}{\mathop{\text{\textit{\textbf{E}}}}}
\newcommand{\ibF}{\mathop{\text{\textit{\textbf{F}}}}}
\newcommand{\ibG}{\mathop{\text{\textit{\textbf{G}}}}}
\newcommand{\ibH}{\mathop{\text{\textit{\textbf{H}}}}}
\newcommand{\ibI}{\mathop{\text{\textit{\textbf{I}}}}}
\newcommand{\ibJ}{\mathop{\text{\textit{\textbf{J}}}}}
\newcommand{\ibK}{\mathop{\text{\textit{\textbf{K}}}}}
\newcommand{\ibL}{\mathop{\text{\textit{\textbf{L}}}}}
\newcommand{\ibM}{\mathop{\text{\textit{\textbf{M}}}}}
\newcommand{\ibN}{\mathop{\text{\textit{\textbf{N}}}}}
\newcommand{\ibO}{\mathop{\text{\textit{\textbf{O}}}}}
\newcommand{\ibP}{\mathop{\text{\textit{\textbf{P}}}}}
\newcommand{\ibQ}{\mathop{\text{\textit{\textbf{Q}}}}}
\newcommand{\ibR}{\mathop{\text{\textit{\textbf{R}}}}}
\newcommand{\ibS}{\mathop{\text{\textit{\textbf{S}}}}}
\newcommand{\ibT}{\mathop{\text{\textit{\textbf{T}}}}}
\newcommand{\ibU}{\mathop{\text{\textit{\textbf{U}}}}}
\newcommand{\ibV}{\mathop{\text{\textit{\textbf{V}}}}}
\newcommand{\ibW}{\mathop{\text{\textit{\textbf{W}}}}}
\newcommand{\ibX}{\mathop{\text{\textit{\textbf{X}}}}}
\newcommand{\ibY}{\mathop{\text{\textit{\textbf{Y}}}}}
\newcommand{\ibZ}{\mathop{\text{\textit{\textbf{Z}}}}}

\newcommand{\ibx}{\mathop{\text{\textit{\textbf{x}}}}}

%\newcommand{\Comp}{\mathop{\mathrm{C}}\nolimits}
%\newcommand{\Komp}{\mathop{\mathrm{K}}\nolimits}
%\newcommand{\Domp}{\mathop{\mathsf{D}}\nolimits}%複体のホモトピー圏
%\newcommand{\Comp}{\mathrm{C}}
%\newcommand{\Komp}{\mathrm{K}}
%\newcommand{\Domp}{\mathsf{D}}%複体のホモトピー圏

\newcommand{\Comp}{\mathop{\mathrm{C}}\nolimits}
\newcommand{\Komp}{\mathop{\mathsf{K}}\nolimits}
\newcommand{\Domp}{\mathord{\mathsf{D}}}%\nolimits}
\newcommand{\Kompl}{\mathop{\mathsf{K}^\mathrm{+}}\nolimits}
\newcommand{\Kompu}{\mathop{\mathsf{K}^\mathrm{-}}\nolimits}
\newcommand{\Kompb}{\mathop{\mathsf{K}^\mathrm{b}}\nolimits}
\newcommand{\Dompl}{\mathop{\mathsf{D}^\mathrm{+}}\nolimits}
\newcommand{\Dompu}{\mathop{\mathsf{D}^\mathrm{-}}\nolimits}
\newcommand{\Dompb}{\mathop{\mathsf{D}^\mathrm{b}}\nolimits}
\newcommand{\Dompbf}{\mathop{\mathsf{D}_\mathrm{f}^\mathrm{b}}\nolimits}
\newcommand{\Domplb}{\mathop{\mathsf{D}^\mathrm{lb}}\nolimits}




\newcommand{\CCat}{\Comp(\cat)}
\newcommand{\KCat}{\Komp(\cat)}
\newcommand{\DCat}{\Domp(\cat)}%圏Cの複体のホモトピー圏
\newcommand{\HOM}{\mathop{\mathscr{H}\hspace{-2pt}om}\nolimits}%内部Hom
\newcommand{\RHOM}{\mathop{\mathrm{R}\hspace{-1.5pt}\HOM}\nolimits}

%\newcommand{\muS}{\mathop{\mathrm{SS}}\nolimits}
\newcommand{\muS}[1][]{\mathop{\mathrm{SS}}\nolimits_{#1}}
\newcommand{\RG}{\mathop{\mathrm{R}\hspace{-0pt}\Gamma}\nolimits}
\newcommand{\RHom}{\mathop{\mathrm{R}\hspace{-1.5pt}\Hom}\nolimits}
%\newcommand{\Rder}{\mathrm{R}}
\newcommand{\Rder}[1][]{\mathop{\mathrm{R}^{#1}\hspace{-2pt}}\nolimits}

\newcommand{\simar}{\mathrel{\overset{\sim}{\rightarrow}}}%同型右矢印
\newcommand{\simarr}{\mathrel{\overset{\sim}{\longrightarrow}}}%同型右矢印
\newcommand{\simra}{\mathrel{\overset{\sim}{\leftarrow}}}%同型左矢印
\newcommand{\simrra}{\mathrel{\overset{\sim}{\longleftarrow}}}%同型左矢印

\newcommand{\hocolim}{{\mathrm{hocolim}}}
\newcommand{\indlim}[1][]{\mathop{\varinjlim}\limits_{#1}}
\newcommand{\sindlim}[1][]{\smash{\mathop{\varinjlim}\limits_{#1}}\,}
\newcommand{\Pro}{\mathrm{Pro}}
\newcommand{\Ind}{\mathrm{Ind}}
\newcommand{\prolim}[1][]{\mathop{\varprojlim}\limits_{#1}}
\newcommand{\sprolim}[1][]{\smash{\mathop{\varprojlim}\limits_{#1}}\,}

\newcommand{\Sh}{\mathrm{Sh}}
\newcommand{\PSh}{\mathrm{PSh}}

\newcommand{\rmD}{\mathrm{D}}

\newcommand{\Lloc}[1][]{\mathord{\mathcal{L}^1_{\mathrm{loc},{#1}}}}
\newcommand{\ori}{\mathord{\mathrm{or}}}
\newcommand{\Db}{\mathord{\mathscr{D}b}}

\newcommand{\codim}{\mathop{\mathrm{codim}}\nolimits}



\newcommand{\gld}{\mathop{\mathrm{gld}}\nolimits}
\newcommand{\wgld}{\mathop{\mathrm{wgld}}\nolimits}


\newcommand{\tens}[1][]{\mathbin{\otimes_{\raise1.5ex\hbox to-.1em{}{#1}}}}
\newcommand{\etens}{\mathbin{\boxtimes}}
\newcommand{\ltens}[1][]{\mathbin{\overset{\mathrm{L}}\tens}_{#1}}
\newcommand{\mtens}[1][]{\mathbin{\overset{\mathrm{\mu}}\tens}_{#1}}
\newcommand{\lltens}[1][]{{\mathop{\tens}\limits^{\mathrm{L}}_{#1}}}
\newcommand{\letens}{\overset{\mathrm{L}}{\etens}}
\newcommand{\detens}{\underline{\etens}}
\newcommand{\ldetens}{\overset{\mathrm{L}}{\underline{\etens}}}
\newcommand{\dtens}[1][]{{\overset{\mathrm{L}}{\underline{\otimes}}}_{#1}}

\newcommand{\blk}{\mathord{\ \cdot\ }}
\newcommand{\mres}[2][]{{\left.{#1}\right\rvert}_{#2}}


%\newcommand{\hocolim}{{\mathrm{hocolim}}}
%\newcommand{\indlim}[1][]{\mathop{\varinjlim}\limits_{#1}}
%\newcommand{\sindlim}[1][]{\smash{\mathop{\varinjlim}\limits_{#1}}\,}
%\newcommand{\Pro}{\mathrm{Pro}}
%\newcommand{\Ind}{\mathrm{Ind}}
%\newcommand{\prolim}[1][]{\mathop{\varprojlim}\limits_{#1}}
%\newcommand{\sprolim}[1][]{\smash{\mathop{\varprojlim}\limits_{#1}}\,}
\newcommand{\proolim}[1][]{\mathop{\text{\rm``{$\varprojlim$}''}}\limits_{#1}}
\newcommand{\sproolim}[1][]{\smash{\mathop{\rm``{\varprojlim}''}\limits_{#1}}}
\newcommand{\inddlim}[1][]{\mathop{\text{\rm``{$\varinjlim$}''}}\limits_{#1}}
\newcommand{\sinddlim}[1][]{\smash{\mathop{\text{\rm``{$\varinjlim$}''}}\limits_{#1}}\,}
\newcommand{\ooplus}{\mathop{\text{\rm``{$\oplus$}''}}\limits}
\newcommand{\bbigsqcup}{\mathop{``\bigsqcup''}\limits}
\newcommand{\bsqcup}{\mathop{``\sqcup''}\limits}
\newcommand{\dsum}[1][]{\mathbin{\oplus_{#1}}}

\newcommand{\Fct}{\mathop{\mathsf{Fct}}\nolimits}

\newcommand{\pb}[1][]{\mathbin{\mathop{\times}\limits_{#1}}}
\newcommand{\dpb}[1][]{\mathbin{\mathop{\times}\limits_{#1}}}
\newcommand{\tpb}[1][]{\mathbin{\mathop{\times}\nolimits_{#1}}}





%================================================
% 自前の定理環境
%   https://mathlandscape.com/latex-amsthm/
% を参考にした
\newtheoremstyle{mystyle}%   % スタイル名
    {5pt}%                   % 上部スペース
    {5pt}%                   % 下部スペース
    {}%              % 本文フォント
    {}%                  % 1行目のインデント量
    {\bfseries}%                      % 見出しフォント
    {.}%                     % 見出し後の句読点
    {12pt}%                     % 見出し後のスペース
    {\thmname{#1}\thmnumber{ #2}\thmnote{{\hspace{2pt}\normalfont (#3)}}}% % 見出しの書式

\theoremstyle{mystyle}
\newtheorem{AXM}{公理}[section]
\newtheorem{DFN}[AXM]{定義}
\newtheorem{THM}[AXM]{定理}
\newtheorem*{THM*}{定理}
\newtheorem{PRP}[AXM]{命題}
\newtheorem{LMM}[AXM]{補題}
\newtheorem{CRL}[AXM]{系}
\newtheorem{EG}[AXM]{例}
\newtheorem*{EG*}{例}
\newtheorem{RMK}[AXM]{注意}
\newtheorem{CNV}[AXM]{約束}
\newtheorem{CMT}[AXM]{コメント}
\newtheorem*{CMT*}{コメント}
\newtheorem{NTN}[AXM]{記号}
\newtheorem{PRB}[AXM]{問題}
\newtheorem{OBS}[AXM]{観察}
\newtheorem{CNJ}[AXM]{予想}

% 定理環境ここまで
%====================================================

%\usepackage{framed}
\definecolor{lightgray}{rgb}{0.75,0.75,0.75}
%\renewenvironment{leftbar}{%
%  \def\FrameCommand{\textcolor{lightgray}{\vrule width 4pt} \hspace{10pt}}% 
%  \MakeFramed {\advance\hsize-\width \FrameRestore}}%
%{\endMakeFramed}
%\newenvironment{redleftbar}{%
%  \def\FrameCommand{\textcolor{lightgray}{\vrule width 1pt} \hspace{10pt}}% 
%  \MakeFramed {\advance\hsize-\width \FrameRestore}}%
% {\endMakeFramed}


% tcolorbox パッケージの読み込み（高度な装飾とページ跨ぎのライブラリを追加）
\usepackage{tcolorbox}
\tcbuselibrary{skins, breakable}

\definecolor{lightgray}{rgb}{0.75,0.75,0.75}

% DFN環境を直接tcolorboxで装飾する（leftbarの再現）
\tcolorboxenvironment{DFN}{
  blanker,                              % 背景色やデフォルトの枠線を透明にする
  breakable,                            % ページ跨ぎを許可する（framedの代わり）
  left=14pt,                            % 左側の余白（線の太さ4pt + 隙間10pt）
  borderline west={4pt}{0pt}{lightgray},% 左側(west)に太さ4ptのlightgrayの線を引く
  before skip=10pt,                     % 環境の上の余白
  after skip=10pt                       % 環境の下の余白
}

% DFN環境を直接tcolorboxで装飾する（leftbarの再現）
\tcolorboxenvironment{PRP}{
  enhanced,                             % 高度な装飾を有効化（skinsライブラリが必要）
  breakable,                            % 箱のページまたぎを許す
  boxrule=0pt,                          % デフォルトの枠線をすべて消す（emptyの代わり）
  frame hidden,                         % デフォルトの枠を非表示にする
  sharp corners,                        % 角を直角にする（デフォルトの丸みを消す）
  colback=black!5,                      % 背景色を塗る（\fillの代わり）
  borderline west={2pt}{0pt}{black},    % 左側に2ptの黒線を引く（\drawの代わり）
  left=12pt,                            % 左側の余白（線の太さ2pt + 隙間10pt）
  before skip=10pt,                     % 環境の上の余白
  after skip=10pt                       % 環境の下の余白
}
% DFN環境を直接tcolorboxで装飾する（leftbarの再現）
\tcolorboxenvironment{THM}{
  enhanced,                             % 高度な装飾を有効化（skinsライブラリが必要）
  breakable,                            % 箱のページまたぎを許す
  boxrule=0pt,                          % デフォルトの枠線をすべて消す（emptyの代わり）
  frame hidden,                         % デフォルトの枠を非表示にする
  sharp corners,                        % 角を直角にする（デフォルトの丸みを消す）
  colback=black!5,                      % 背景色を塗る（\fillの代わり）
  borderline west={2pt}{0pt}{black},    % 左側に2ptの黒線を引く（\drawの代わり）
  left=12pt,                            % 左側の余白（線の太さ2pt + 隙間10pt）
  before skip=10pt,                     % 環境の上の余白
  after skip=10pt                       % 環境の下の余白
}

% 参考：例(EG)環境に redleftbar（細い線）を適用する場合
\tcolorboxenvironment{EG}{
  blanker,
  breakable,
  left=11pt,                            % 左側の余白（線の太さ1pt + 隙間10pt）
  borderline west={1pt}{0pt}{lightgray},% 左側に太さ1ptの線を引く
  before skip=10pt,
  after skip=10pt
}


\renewcommand{\Re}{\mathop{\mathrm{Re}}\nolimits}


% =================================





% ---------------------------
% new definition macro
% ---------------------------
% 便利なマクロ集です

% 内積のマクロ
%   例えば \inner<\pphi | \psi> のように用いる
\def\inner<#1>{\langle #1 \rangle}

% ================================
% マクロを追加する場合のスペース 

%=================================



% 位相空間
\newcommand{\Int}[1][]{\mathop{\mathrm{Int}}\nolimits_{#1}}
\newcommand{\Cl}[1][]{\mathop{\mathrm{Cl}}\nolimits_{#1}}
\newcommand{\Bd}[1][]{\mathop{\mathrm{Bd}}\nolimits_{#1}}
\newcommand{\bd}{\mathord{\partial}}
\newcommand{\Fr}[1][]{\mathop{\mathrm{Fr}}\nolimits_{#1}}


\newcommand{\mtan}[1][]{T\hspace{-1.75pt}{#1}}
\newcommand{\mtp}[1][]{{}^{t\!}{#1}} % transpose of the morphism
%\newcommand{\mpb}[1][]{{}^{t\!}{#1}'} % pull-back of the morphism
\newcommand{\mpb}[1][]{{#1}_{\pi}} % pull-back of the morphism
\newcommand{\mdiff}[1][]{{#1}_{d}} % pull-back of the morphism
\newcommand{\mptan}[1][]{{#1}_{\tau}} % pull-back of the morphism


%\newcommand{\pb}[1][]{\mathbin{\mathop{\times}\limits_{#1}}}

\newcommand{\cotb}[1][]{T^{\ast}{#1}}
\newcommand{\conm}[2][]{T_{#1}^{\hspace{0.7pt}\ast}{#2}}
\newcommand{\norb}[2][]{T_{#1}{#2}}

\newcommand{\cotz}[1][]{T^{\ast}_{#1}{#1}}

\newcommand{\cotn}[1][]{\dot{T}^{\ast}{#1}}
\newcommand{\conmn}[2][]{\dot{T}_{#1}^{\hspace{0.7pt}\ast}{#2}}


\newcommand{\muhom}{\mu hom}
\newcommand{\muHom}[1][]{\mathrm{Hom}^\mu_{\raise1.5ex\hbox to.1em{}#1}}

\newcommand{\mucat}[2][]{\mathsf{D}^{\mathrm{b}}_{#1}(#2)}
%\newcommand{\mucat}[2][]{\mathsf{D}^{\mathrm{b}}_{#1}(#2)}

\newcommand{\conv}[1][]{\mathbin{\mathop{\circ}\limits_{#1}}}


\newcommand{\lopen}{\mathopen{]}}
\newcommand{\lclosed}{\mathopen{[}}
\newcommand{\ropen}{\mathclose{[}}
\newcommand{\rclosed}{\mathopen{]}}

% 4種類の区間をすべて定義
\DeclarePairedDelimiter{\openintvl}{]}{[}     % ]a, b[
\DeclarePairedDelimiter{\closedintvl}{[}{]}       % [a, b]
\DeclarePairedDelimiter{\lopenrclosedintvl}{]}{]}    % ]a, b] (左半開)
\DeclarePairedDelimiter{\lcosedropenintvl}{[}{[}    % [a, b[ (右半開)


\newcommand{\sm}{\raisebox{2.33pt}{~\rule{6.4pt}{1.3pt}~}}
\newcommand{\smi}{\smallsetminus}

\newcommand{\smid}{\mathrel{;}}

\newcommand{\hfun}[1][]{\mathcal{B}_{#1}}
\newcommand{\mfun}[1][]{\mathcal{C}_{#1}}
\newcommand{\mildhfun}[1][]{\widehat{\mathcal{B}}_{#1}}
\newcommand{\mildmfun}[1][]{\widehat{\mathcal{C}}_{#1}}

\newcommand{\micro}[1]{\mu_{#1}\hspace{0pt}}

\newcommand{\rrp}{\rr^{+}}
\newcommand{\Dbconic}{\mathop{\mathsf{D}^\mathrm{b}_{\rrp}}\nolimits}
\newcommand{\iu}{\sqrt{-1}}

\DeclareMathOperator{\WFa}{\mathrm{WF}_\mathrm{A}} % 解析的波面集合
\DeclareMathOperator{\WF}{\mathrm{WF}} % 波面集合
\DeclareMathOperator{\SingSpec}{\mathrm{S}\text{.}\mathrm{S}\text{.}} % 特異性スペクトル
\DeclareMathOperator{\SingSupp}{\mathrm{singsupp}} % 特異台
\newcommand{\Char}{\mathop{\mathrm{Char}}\nolimits}

% ----------------------------
% documenet 
% ----------------------------
% 以下, 本文の執筆スペースです. 
% Your main code must be written between 
% begin document and end document.
% ---------------------------

\title{手を動かして学ぶ\footnote{「まなぶ」ではないことに注意．} 代数解析【佐藤超関数編】}
\author{Toshihiro Oshiba}
\date{\today}
\begin{document}
\maketitle

\section*{はじめに}
2026/06/26に実施された北大数学オムニバス発表会のノート．
\section{複素関数論}
\(D\)を複素平面\(\cc\)の開集合，
%\(z\in D\)を\(D\)の点とする．
\(f\)を\(D\)上の\(C^1\)級関数とする．すなわち
\(z=x+\iu y\)とするとき，\(f(x,y)\)が連続微分可能であるとする．
\(f\)が\(z\in D\)において正則であるとは，
\[
    \bar{\partial}f(z)=0
\]
がなりたつことをいう．
\begin{itemize}
  \item \(\cO(D)\coloneqq\left\{\text{\(D\)上の正則関数}\right\}\)とおく．
  \item 正則関数の実軸への制限として表される関数を実解析関数とよぶ．
  \item \(\cA(D)\coloneqq\left\{\text{\(D\)上の実解析関数}\right\}\)とおく．
\end{itemize}

\begin{THM}[コーシーの積分定理]
  \(D\)を\(\cc\)内の単連結な領域，例えば\(D^2\)とする．
  \(f\in\cO(D)\)とし，
  \(C\)を\(D\)内の区分的に\(C^1\)級な単純閉曲線，
  例えば半径\(0<r<1\)の円周とする．
  このとき，
  \[
    \int_{C}f(\zeta)d\zeta=0
  \]
  である．
\end{THM}
\begin{THM}[コーシーの積分定理]
  \(D\)を\(\cc\)内の単連結な領域，例えば\(D^2\)とする．
  \(f\in\cO(D)\)とし，
  \(C\)を\(D\)内の区分的に\(C^1\)級な単純閉曲線，
  例えば半径\(0<r<1\)の円周とする．
  このとき，\(C\)に囲まれる点\(z_0\)に対し，
  \[
    f(z_0)=\frac{1}{2\pi\sqrt{-1}}\int_{C}\frac{1}{\zeta-z_0}f(\zeta)d\zeta=0
  \]
  がなりたつ．
\end{THM}

\(f\)が\(z_0\)で正則であることと，\(f\)が\(z_0\)の近傍でテイラー展開
\[
  f(z)=\sum_{n=0}^{\infty}a_n(z-z_0)^n
\]
を持つことは同値である．

\begin{DFN}
  \(D_{1}\subset D_{2}\)を\(\cc\)内の領域とする．
  \(f_{1}\in D_{1}\)，
  \(f_{2}\in D_{2}\)とする．\(f_{1}=\mres[f_{2}]{D_{1}}\)のとき，
  \(f_{2}\)は\(f_{1}\)の\(D_{2}\)への解析接続であるという．
\end{DFN}
\begin{THM}[一致の定理]
  \(D_{1}\subset D_{2}\)を\(\cc\)内の開集合とする．
  制限射\(\cO(D_2)\to \cO(D_1)\)は単射である．
\end{THM}

\section{境界値問題からの動機}
2次元開円盤を\(D\)で表す．すなわち
\[
    D\coloneqq\left\{z\in\cc\smid \left|z\right|<1\right\}
\]である．\(D\)の境界である1次元円周を\(S^1\)で表す．
すなわち\[
    S^1\coloneqq \bd{D}=\left\{z\in\cc\smid \left|z\right|=1\right\}
\]である．次の問題を考える．
\begin{PRB}[境界値問題（ディリクレ問題）]
    \(f\)を\(S^1\)上の\(C^\infty\)級関数とする．
    このとき\(f\)を境界値とする\(D\)上の調和関数\(u\)を求めよ．
    すなわち，次をみたす関数\(u\)を求めよ．
    \[
      \begin{cases*}
        \Delta u=0&on \(D\),\\
        \mres[u]{S^1}=f.
      \end{cases*}
    \]
\end{PRB}

この問題の解法のひとつにPoisson積分によるものがある．
\(f\in C^\infty(S^{1})\)に対し，\(D\)上の関数\(\cP{f}\)を
\[
  \cP{f}(z)\coloneqq
  \frac{1}{2\pi}\int_{0}^{2\pi}\frac{1-\left|z\right|^2}{\left|e^{\sqrt{-1}\theta}-z\right|}f(\theta) d\theta
\]
で定める．\(\cP{f}\)を\(f\)のポアソン積分とよぶ．
\[
  K_{z}=\frac{1-\left|z\right|^2}{\left|e^{\sqrt{-1}\theta}-z\right|}
\]
とおいて，\(K_z\)をポアソン核とよぶこともある．
\(\cP{f}\)は\(D\)上の調和関数である．

\(\cP\)は線型写像
\[
  \cP\colon C^{\infty}(S^{1})\to \cH(D)
\]
を定める．ただし
\[
  \cH(D)\coloneqq \left\{\text{\(D\)上の調和関数}\right\}
\]
である．

\begin{OBS}
  \(f,g\in C^{\infty}(S^1)\)に対し\(\cP{f}=\cP{g}\)とすると，
  境界値の定義により，\(f=\mres[\cP{f}]{S^{1}}=\mres[\cP{g}]{S^{1}}=g\)である．
  したがって，\(\cP\)は単射である．

  関数\(\dfrac{1}{1-z}\)を考える．
  \(\dfrac{1}{1-z}\in \cH(D)\)である．
  \(\mres[\dfrac{1}{1-z}]{S^{1}}\)は\(1\in S^1\)で微分できずポアソン積分で表せない．
  すなわち，\(\cP\)は全射ではない．
\end{OBS}

以上の観察から次のような問題が考えられる．
\begin{PRB}
  \(\cP\)が同型となるような ``関数'' のクラスを決定せよ．
\end{PRB}
\section{シュヴァルツ分布}
\subsection{定義}

\(n\in\nn\)を自然数とする．
\(f\in C^{\infty}(S^{1})\)に対し
\[
  \nu_{n}(f)\coloneqq
  \sum_{k=0}^{n}\max_{0\leqq\theta<2\pi}\left|\frac{d^k}{d\theta^k}f(\theta)\right|
\]
とおく．このとき，
\[
  d(f,g)=\sum_{n=0}^{\infty}\frac{1}{2^n}\frac{\nu_{n}(f-g)}{1+\nu_{n}(f-g)}
\]
は\(C^{\infty}(S^1)\)上の距離を定める．
この距離で\(C^{\infty}(S^1)\)を線形位相空間とみなす．

連続線型写像\(T\colon C^{\infty}(S^1)\to \cc\)を
\emph{シュヴァルツ分布} (Schwartz distribution) とよぶ．
\[
  \Db(S^1)\coloneqq \left\{\text{\(S^1\)上のシュヴァルツ分布}\right\}
\]
とおく．

\begin{RMK}
  実数直線\(\rr\)上の分布は\(C^{\infty}_{c}(\rr)\)の
  （位相的な）双対空間
  の元として定める．
\end{RMK}

分布\(T\)の\(f\)での値\(T(f)\)をペアリング
\[
  \inner<T,f>
\]
で表すこともある．
\subsection{例や演算・性質}

\begin{EG}
  ディラックのデルタ関数を
  \[
    \delta(f)=f(0)
  \]
  で定める．これは\(S^1\)上の分布である．
\end{EG}
\begin{EG}
  実数直線上のディラックのデルタ関数を
  \[
    \delta(f)=f(0)
  \]
  で定める．これは\(\rr\)上の分布である．

  ところで，ペアリングの表示から
  \[
    \inner<\delta,f>=\int_{-\infty}^{\infty}\delta(x)f(x)dx
  \]
  ともかける．
\end{EG}

分布の立場では微分を引数（テスト関数）の微分に押し付けて定義する．
\begin{DFN}
  \(T\in\Db(\rr)\)とする．
  \(T\)の微分を
  \[
    T'(f)=-T(f')=-\inner<T,f'>
  \]
  で定める．
\end{DFN}
\begin{EG}
    ヘヴィサイド関数\(Y\)を
  \[
    \inner<Y,f>\coloneqq
    \int_{0}^{\infty}f(x)dx
  \]
  で定める．\(Y(x)\)の微分を求める．
  \begin{align*}
    \inner<Y',f>
    &=-\inner<Y,f'>\\
    &=-\int_{0}^{\infty}f'(x)dx\\
    &=-\left[f(x)\right]_{0}^{\infty}\\
    &=f(0)-0=\delta(f).
  \end{align*}
  よって，\(Y'=\delta\)である．
\end{EG}
\(C^\infty\)級関数は分布とみなせる．
\begin{DFN}
  \(C^\infty\)級関数\(f\in C^\infty(S^1)\)に対し，
  \(T_{f}\)を
  \[
    T_{f}(g)=\frac{1}{2\pi}\int_{0}^{2\pi}f(\theta)f(\theta)d\theta
  \]
  で定める．\(T_{f}\)は\(S^1\)上の分布である．
\end{DFN}

\begin{PRP}
  \(T_{\blk}\colon C^{\infty}(S^1)\hookrightarrow \Db(S^1)\)は連続な線形単射である．
\end{PRP}
\subsection{問題について}

\begin{OBS}
  \(T\)に対し，
  \[
    \tilde{\cP}(T)=T(K_{z})=\inner<T,K_{z}>
  \]
  とおく．ただし，\(K_z\)はポアソン核
  である．
  このとき，\(\tilde{\cP}(T)\)は\(D\)上の調和関数になる．
  また，\(f\in C^{\infty}\)に対し
  \[
    \tilde{\cP}(T_{f})=\frac{1}{2\pi}\int_{0}^{2\pi}\frac{1-\left|z\right|^2}{\left|e^{\sqrt{-1}\theta}-z\right|}f(\theta) d\theta=\cP{f}
  \]
  であり，\(\tilde{\cP}\)は\(\cP\)の\(\Db(S^1)\)への拡張になっている．
  \(\tilde{\cP}\)も\(\cP\)で表す．

  \(T\in \Db(S^1)\)とする．\(\cP(T)=0\)とする．
  このとき，フーリエ級数が全て0になることが示せて，\(T=0\)となる．
  よって\(\cP\)は単射である．

  \(e^{\frac{1}{1-z}}\)は調和関数である．
  \(\cP(T)=e^{\frac{1}{1-z}}\)となる\(T\in \Db(S^1)\)が
  存在したと仮定する．
  これをフーリエ級数展開すると
  \(T\notin \Db(S^1)\)であることが示せる．
  したがって，\(\cP\colon \Db(S^1)\to \cH(D)\)は全射でない．
\end{OBS}


\section{佐藤超関数}

\subsection{デルタ関数再論}
実数直線上のディラックのデルタ関数の定義を思い出す．
\[
  \inner<\delta,f>=\int_{-\infty}^{\infty}\delta(x)f(x)dx
\]
だった．これは原点周りの円周\(C\)をとると
コーシーの積分公式
\[
  f(0)=\frac{1}{2\pi\sqrt{-1}}\int_{C}\frac{1}{z}f(z)dz
\]
と似ている．
そこで次のようにしてみる．
\(C\)の上を\(C_+\)，下を\(C_{-}\)とすると，
\begin{align*}
  \frac{1}{2\pi\sqrt{-1}}\int_{C}\frac{1}{z}f(z)dz
  &=
  \frac{1}{2\pi\sqrt{-1}}\int_{C_{+}}\frac{1}{z}f(z)dz
  +\frac{1}{2\pi\sqrt{-1}}\int_{C_{-}}\frac{1}{z}f(z)dz\\
  &=
  \frac{1}{2\pi\sqrt{-1}}\int_{-C_{+}}-\frac{1}{z}f(z)dz
  -\frac{1}{2\pi\sqrt{-1}}\int_{C_{-}}-\frac{1}{z}f(z)dz\\
  &=
  \frac{1}{2\pi\sqrt{-1}}\left(
  \int_{-1+\sqrt{-1}0}^{+1+\sqrt{-1}0}-\frac{1}{z}f(z)dz
  -\int_{-1-\sqrt{-1}0}^{+1-\sqrt{-1}0}-\frac{1}{z}f(z)dz\right)\\
  &=
  \int_{-1}^{+1}\left(-\frac{1}{2\pi\sqrt{-1}}\right)\left(
  \frac{1}{x+\sqrt{-1}0}
  -\frac{1}{x-\sqrt{-1}0}\right)f(x)dx
\end{align*}
デルタ関数の式と比較すれば
\[
  \delta(x)=\left(-\frac{1}{2\pi\sqrt{-1}}\right)\left(
  \frac{1}{x+\sqrt{-1}0}
  -\frac{1}{x-\sqrt{-1}0}\right)
\]
と解釈できる．
\subsection{定義}
\(I\)を実数直線内の開区間とする．
\(\varOmega\)を\(I\)を閉集合として含む複素平面の開集合とする．
このとき，\(I\)上の佐藤超関数の空間\(\cB(I)\)を
\[
    \cB(I)\coloneqq \cO(\varOmega\sm I)/\cO(\varOmega)
\]で定める．\(\cB(I)\)は\(\varOmega\)の取り方によらない．
（あるいは余極限\[
    \cB(I)=\colim_{\varOmega} \cO(\varOmega\sm I)/\cO(\varOmega)
\]として定めてもよい．）
\(a<b\)を用いて\(I=(a,b)\)とかくとき，
\[
    \cO(\varOmega \sm I)
    \cong 
    \cO(\varOmega\cap H^{+})
    \oplus
    \cO(\varOmega\cap H^{-})
\]が成り立つので，\(\cB(I)\)の元\(f\)に対し，
\((F_{+},F_{-})\in\cO(\varOmega\cap H^{+})
    \oplus
    \cO(\varOmega\cap H^{-})
\)を用いて，\[
    f=[(F_{+},F_{-})]
\]とかける．これを\[
    f(x)=F_{+}(x+\iu 0)-F_{-}(x-\iu 0)
\]ともかく．
\subsection{例や演算・性質}

\begin{EG}
  ディラックのデルタ関数を
  \[
    \delta(x)=\left[-\frac{1}{2\pi\sqrt{-1}}\frac{1}{z}\right]
  \]
  で定める．
\end{EG}

微分は定義関数の微分の類として定める．
\begin{EG}
  ヘヴィサイド関数\(Y\)を
  \[
    Y(x)=\left[-\frac{1}{2\pi\sqrt{-1}}\log(-z)\right]
  \]
  で定める．\(Y(x)\)の微分を求める．
  \begin{align*}
    Y'(x)
    &=\left[-\frac{1}{2\pi\sqrt{-1}}\log(-z)\right]\\
    &=\left[-\frac{1}{2\pi\sqrt{-1}}\frac{1}{-z}\times(-1)\right]\\
    &=\left[-\frac{1}{2\pi\sqrt{-1}}\frac{1}{z}\right]\\
    &=\delta(x)
  \end{align*}
  となる．
\end{EG}

\begin{EG}
  \((x+\sqrt{-1}0)^\alpha\)を
  \[
    (x+\sqrt{-1}0)^\alpha\coloneqq
    \left[z^{\alpha},0\right]
  \]
  で定める．
\end{EG}

\begin{EG}
  \(f\)を実解析関数とする．このとき
  \[
    [(f,0)]
  \]
  とすると超関数が定まる．よって\(\cA\hookrightarrow\cB\)が定まる．
\end{EG}


\subsection{問題について}

コンパクトな実解析多様体，例えば\(S^1\)上では，
\(\cB(S^1)\)は\(\cA(S^1)\)の双対空間と
同型であることが知られている．
\begin{THM}
  \(\cP\)は\[
    \cP\colon \cB(S^1)\to \cH(D)
  \]
  に拡張される．さらにこれは同型である．
\end{THM}
\section{応用：Helgason予想}
\(C^{\infty}(S^1)\)に対し，
\[
  \cP_{s}f(z)\coloneqq \frac{1}{2\pi}\int_{0}^{2\pi}K_{z}^{s}f(\theta)d\theta
\]
とおく．
\[
  \cH_{\lambda}(D)
  \coloneqq
  \left\{
    u\in C^{\infty}(D)\smid
    (1-x^2-y^2)^2\Delta{u}(z)=\lambda u(z)
  \right\}
\]
とするとき，線型写像
\[
  \cP_{s}\colon C^{\infty}(S^1)\to
  \cH_{\lambda}(D)
\]
が定まる．

Helgasonは次を示した．
\begin{THM}[Helgason, 1970]
  \(\cP_{s}\)は\[
    \cP_{s}\colon \cB(S^1)\to \cH_{4s(s-1)}(D)
  \]
  に拡張される．さらにこれは同型である．
\end{THM}

Helgasonはさらに次の予想を提示した．
\begin{CNJ}[Helgason予想]
  一般のリーマン対称空間上の``球関数''は，
  ``境界''上の佐藤超関数のポアソン積分を用いて表せるであろう．
\end{CNJ}

\begin{THM}[Kashiwara--Kowata--Minemura--Okamoto--Oshima--Tanaka]
  Helgason予想は正しい．
\end{THM}
この予想の解決には超局所解析の理論が用いられた．

\section{超局所解析}

関数の特異性を方向別に調べる．
まず
\[
  \SingSupp(f)\coloneqq\left\{
    x\in\varOmega\smid 
    \text{\(f\)は\(x\)で解析的でない}
  \right\}
\]
とおく．

\begin{DFN}
  \(f=[F_{+},F_{-}]\in\cB(\varOmega)\)を
  超関数とする．
  \((x;\xi)\in\sqrt{-1}\cotn[\varOmega]\)とする．
  \(f\)が\((x;\xi)\)で\emph{超局所解析的} (microanalytic) で
  あるとは，実数\(\varepsilon\)で，\(\xi>0\)なら
  \(F_{+}\)が，\(\xi<0\)なら\(F_{-}\)が，
  \((x;-\varepsilon\xi)\)の近傍に解析接続できるものが存在する
  ことをいう．

  \(f\)の\emph{特異性スペクトル} (singularity 
  spectrum) 
  \(\SingSpec(f)\)あるいは
  \emph{波面集合} (wavefront set) \(\WF(f)\)を
  \[
    \SingSpec(f)\coloneqq
    \left\{
      (x;\xi)\in\sqrt{-1}\cotn[\varOmega]
      \smid
      \text{\(f\)は\((x;\xi)\)で超局所解析的でない}
    \right\}\cup \supp(f)
  \]
  と定義する．
\end{DFN}

\(\pi\colon \cotb[\varOmega]\to \varOmega\)を
射影\(\pi(x,\xi)=x\)とする．
\(\dot{\pi}\)を\(\pi\)の\(\cotn[\varOmega]\)への制限とする．
このとき，
\[
  \dot{\pi}(\SingSpec(f))=\SingSupp(f)
\]
が成り立つ．

\begin{EG}
  \begin{itemize}
    \item \(\SingSpec{\delta}=\dot{\pi}^{-1}(0)\)
    \item \(\SingSpec{Y}=\dot{\pi}^{-1}(0)\)
    \item \(\SingSpec{\dfrac{1}{x+\sqrt{-1}0}}
    =\left\{(0;\sqrt{-1}\xi)\in\sqrt{-1}\cotn[\rr]\smid
    \xi>0\right\}\)
    \item \(\SingSpec{\dfrac{1}{\inner<x,a>+\sqrt{-1}0}}
    =\left\{(x;\sqrt{-1}\xi)\in\sqrt{-1}\cotn[\rr^n]\smid
    \inner<x,a>,\ \xi\in\rrp a\right\}\)
    \item 
    \begin{align*}
    \SingSpec{
      \dfrac{1}{(x_{1}+\sqrt{-1}0)(x_{2}+\sqrt{-1}0)}
    }
    &=\left\{
      (0;\sqrt{-1}\xi)\in\sqrt{-1}\cotn[\rr^2]\smid
      \xi\geqq0
    \right\}\\
    &\cup 
    \left\{
      (0,x_2;\sqrt{-1}\xi_1,0)
      \in\sqrt{-1}\cotn[\rr^2]\smid
      \xi_1\geqq0
    \right\}\\
    &\cup 
    \left\{
      (x_1,0;0,\sqrt{-1}\xi_2)
      \in\sqrt{-1}\cotn[\rr^2]\smid
      \xi_2\geqq0
    \right\}  
    \end{align*}
  \end{itemize}
\end{EG}

微分作用素に対しても余接束の部分集合を考えることができる．

\begin{DFN}
  \(P\)をを実解析関数を係数とする\(\varOmega\)上の線形編微分作用素とする．
  \(P\)の階数が\(m\)であり，\(P\)が
  \[
    P(x,\partial_x)
    =
    \sum_{\left|\alpha\right|\leqq m}^{}
    a_{\alpha}(x)\partial_{x}^{\alpha}
  \]
  と表されているとする．
  このとき，
  \(P\)の\emph{主要表象} (principal symbol) \(\sigma(P)\)を
  \[
    \sigma(P)(x,\xi)
    =\sum_{\left|\alpha\right|= m}^{}
    a_{\alpha}(x)\xi^{\alpha}
  \]で定める．
  \(P\)の\emph{特性多様体} (characteristic variety) 
  \(\Char(P)\)を
  \[
    \Char(P)
    \coloneqq
    \left\{(x;\xi)\in \sqrt{-1}\cotn[\varOmega]\smid
      \sigma(P)(x,\xi)=0
    \right\}
  \]
  で定める．

  \(\Char{P}=\varnothing\)のとき，\(P\)は\emph{楕円型} (elliptic) であるという．
\end{DFN}

次が主定理．

\begin{THM}[佐藤--H\"ormanderの基本定理]
  \(\varOmega\)を\(\rr^n\)の開集合とする．
  \(P\)を実解析関数を係数とする\(\varOmega\)上の線形編微分作用素とする．
  \(f\)を\(\varOmega\)上の超関数とする．
  \(u\)を微分方程式
  \[
    Pu=f
  \]
  の超関数解とする．このとき次が成り立つ．
  \[
    \SingSpec(u)\subset\Char(P)\cup\SingSpec(f).
  \]
\end{THM}

\begin{EG}
  \(\varOmega=\Int[]D^2\subset\rr^2\)とする．
  \(P=\Delta=\partial_{x}^{2}+\partial_{y}^{2}\)とする．
  微分方程式
  \[
    \Delta u=0
  \]
  を考える．
  \(\sigma(\Delta)=\xi_{1}^2+\xi_{2}^2\)であるから
  \begin{align*}
    \Char{P}= \varnothing.
  \end{align*}
  すなわち\(\Delta\)は楕円型である．
  また，\(\SingSpec0=\varnothing\)であるから，
  \[
    \SingSpec u\subset \varnothing\cup \varnothing =\varnothing
  \]
  したがって，\(\SingSpec {u}=\varnothing\)である．
  底空間に射影して
  \[
    \SingSupp{u}=\varnothing.
  \]
  すなわち，調和関数は実解析的である．
\end{EG}

もっと一般に，次も示された．
\begin{THM}
  楕円型微分方程式
  \(Pu=0\)
  の解は実解析的である．
\end{THM}
これは\emph{楕円型正則性} (elliptic regularity) 定理と呼ばれる，
偏微分方程式論における基本定理である．

\begin{RMK}
  余接束（余法束）上の層\(\cC\)と層の射\(
    \mathrm{sp}\colon \cB\to \dot{\pi}_{\ast}\cC
  \)であって
  \[
    0
    \to 
    \cA
    \to 
    \cB
    \overset{\mathrm{sp}}{\longrightarrow} 
    \dot{\pi}_{\ast}
    \cC
    \to 0
  \]
  が完全となるものがある．
  \(\cC\)の切断を超局所関数とよぶ．
  \(\SingSpec {u}\)は\(\supp \mathrm{sp}(u)\)とみなせる．
\end{RMK}
%\section{a}

\section{おまけ}

多変数の超関数は次のようになる．
\(M\)を\(n\)次元実解析多様体，\(X\)を複素近傍とする．
\[
  \cB_{M}\coloneqq
  \RG_{M}(\cO_{X})\tens[\cc]\ori_{M}[n]
\]
同値な定義として
\begin{equation*}\label{eq:B-def2}
  \cB_{M}\coloneqq
  \RHOM_{\cc_X}(\rmD_{X}'{\cc_X}_{M},\cO_X)
\end{equation*}
というものが得られる．
実際，ポワンカレ・ヴェルディエ双対から，
\[
  \rmD'_{X}(\cc_M)\cong \ori_{M}[-n]
\]
が成り立つ．
これに\(\RHOM_{\cc_{X}}(\blk,\cO_X)\)を当てると
\begin{align*}
  \RHOM_{\cc_X}(\rmD_{X}'{\cc_X}_{M},\cO_X)
  &=\RHOM_{\cc_X}(\ori_{M}[-n],\cO_X)\\
  &=\RHOM_{\cc_X}(\cc_{M}\tens[]\ori_{M}[-n],\cO_X)\\
  &=\RHOM_{\cc_X}(\cc_{M},\cO_X)\tens[]\ori_{M}[n]\\
  &=\RG_{M}\cO_X\tens[]\ori_{M}[n]\\
  &=\cB_M
\end{align*}
である．

実解析函数は
\[
  \cA_{M}=\mres[\cO_{X}]{M}
\]
だが
\[
  \cA_{M}=\cO_{X}\tens[]\cc_{M}
\]と書いてもよい．
このとき，
\(\rmD'_{X}=\RHOM_{\cc_{X}}(\blk,\cc_X)\)とおくと
\[
  \cc_M\cong\rmD'_{X}\rmD'_{X}\cc_{M}
\]
が成り立つから，
\begin{align*}
  \cA_{M}
  &\cong \cO_{X}\tens[]\cc_{M}\\
  &\cong \cO_{X}\tens[]\rmD'_{X}\rmD'_{X}\cc_{M}\\
  &\cong \RHOM_{\cc_{X}}(\rmD'_{X}\cc_{M},\cc_{X})\tens[]\cO_{X}\\
  &\to \RHOM_{\cc_{X}}(\rmD'_{X}\cc_{M},\cc_{X})\tens[]\cO_{X}\\
  &\cong\cB_M
\end{align*}
である．
%===============================================
% 参考文献スペース
%===============================================
\begin{thebibliography}{20} 
  \bibitem[J95]{J95} 神保道夫, 
  {複素関数入門}, 岩波講座 現代数学への入門 4, 岩波書店, 1995; 新装版, 2024. 
  \bibitem[Kan96]{Kan96} 金子晃, 新版 超函数入門, 東京大学出版会, 1996. 
  \bibitem[KKK80]{KKK80} 柏原正樹, 河合隆裕, 木村達雄, 
  代数解析学の基礎, 紀伊國屋数学叢書 18, 紀伊國屋書店, 1980.
  \bibitem[KKKOOT78]{KKKOOT78} M.~Kashiwara, 
  A.~Kowata,K.~Minemura, K.~Okamoto, T.~Oshima, M.~Tanaka, 
  \emph{Eigenfunctions of invariant differential operators on a symmetric space}, Ann. of Math., 107 (1978), pp.1--39.
  \bibitem[KS90]{KS90} Masaki Kashiwara, Pierre Schapira, 
    \textit{Sheaves on Manifolds}, 
    Grundlehren der Mathematischen Wissenschaften, 292, Springer, 1990.
  \bibitem[Mori76]{Mori76} 森本光生, 佐藤超函数入門, 共立出版, 1976. 
  \bibitem[Ok97]{Ok97} 岡本清郷, 
  {フーリエ解析の展望}, すうがくぶっくす 17, 朝倉書店, 1997.
  \bibitem[Sato58]{Sato58} 佐藤幹夫, 超函数の理論, 数学 10 (1958) pp.1--27. 
  \bibitem[Sato69]{Sato69} Mikio Sato, 
  \textit{Hyperfuncitons and 
  Partial Differential equations}, 
  Proc. Int. Conf. on Functional Analysis, Tokyo, pp.91--94, 1969.
  \bibitem[Sato70]{Sato70} Mikio Sato, 
  \textit{Regularity of Hyperfunciton Solutions of 
  Partial Differential equations}, 
  Actes, Congres intern. Math., 
  Tome 2, pp.785--794, 1970.
\end{thebibliography}
%===============================================


\end{document}
